Referensi dalam skripsi ini terdiri dari kumpulan jurnal dan aplikasi serupa yang digunakan sebagai bahan perbandingan. Referensi ini tidak hanya menjelaskan tentang aplikasi yang ada, tetapi juga menguraikan hasil-hasil dari jurnal terkait. Perbedaan dari aplikasi-aplikasi referensi tersebut juga dibahas. Daftar referensi yang digunakan dalam skripsi ini dapat dilihat pada \hyperref[sec:referensi]{\textbf{Lampiran~I}}.

Wahyudi \textit{et al}.~\cite{wahyudi2021pengembangan}, melakukan penelitian terkait pengembangan sistem informasi presensi menggunakan metode \textit{waterfall}. Pada penelitian ini digunakan metode \textit{waterfall} dan menggunakan bahasa pemrograman PHP (\textit{Personal Home Page}) dengan \textit{framework Codeigniter} 3, \textit{framework Bootstrap} sebagai CSS, dan MariaDB sebagai \textit{database}-nya. Aplikasi ini tidak terdapat sistem untuk mendeteksi dan perhitungan jam lembur, tidak terdapat pencatatan lokasi melalui GPS (\textit{Global Positioning System}), dan pengambilan foto ketika melakukan presensi. Sedangkan pada penelitian skripsi aplikasi \textit{presensi online} dikembangkan dengan menggunakan metode \textit{scrum}. Aplikasi penelitian skripsi \textit{online} menggunakan \textit{framework Laravel 9}, \textit{framework Tailwind} sebagai CSS, dan \textit{MySQL} sebagai \textit{database}-nya. Aplikasi penelitian skripsi juga dilengkapi fitur untuk mendeteksi dan menghitung jam lembur, pencatatan lokasi saat presensi melalui GPS (\textit{Global Positioning System}), pengambilan foto saat presensi.

Permana \textit{et al}.~\cite{permana2022aplikasi}, melakukan penelitian terkait aplikasi presensi \textit{online} menggunakan validasi jarak lokasi pengguna berbasis Android pada toko Yonix. Pengembangan aplikasi ini menggunakan metode XP (\textit{Extreme Programming}) dan berbasis Android. Aplikasi ini tidak terdapat pengambilan foto saat presensi dan tidak terdapat sistem untuk mendeteksi jam lembur dan menghitung total jam lembur. Sedangkan aplikasi penelitian skripsi presensi \textit{online} menggunakan metode \textit{scrum} dan berbasis \textit{website}. Aplikasi penelitian skripsi presensi \textit{online} melakukan pengambilan foto saat presensi dan sistem untuk mendeteksi dan menghitung jam lembur.

Ramadhini \textit{et al}.~\cite{ramadhini2023sistem}, melakukan penelitian terkait sistem informasi presensi karyawan berbasis android pada perusahaan asuransi Panin Dai-Ichi Life. Pengembangan aplikasi ini menggunakan metode \textit{waterfall} dan berbasis Android. Aplikasi ini hanya dirancang untuk pegawai saja yang memiliki fitur untuk mengajukan izin atau sakit, tetapi tidak memiliki sistem untuk mendeteksi dan menghitung jam lembur. Sedangkan aplikasi penelitian skripsi presensi \textit{online} dikembangkan dengan metode \textit{scrum} dan berbasis \textit{website}. Aplikasi penelitian skripsi presensi \textit{online} dirancang untuk memiliki 3 \textit{role} yaitu \textit{staff Infrastructure}, \textit{Manager Infrastructure}, dan \textit{Manager Human Resources}. Aplikasi ini tidak memiliki fitur pengajuan izin dan sakit tetapi memiliki sistem untuk mendeteksi dan menghitung jam lembur.

Roosdianto \textit{et al}.~\cite{roosdianto2021rancang}, melakukan penelitian terkait rancang bangun aplikasi sistem informasi absensi karyawan \textit{online}. Pegembangan aplikasi ini menggunakan metode \textit{waterfall} dan menggunakan bahasa pemrograman PHP (\textit{Personal Home Page}) dengan \textit{framework CodeIgniter} dan \textit{framework Bootstrap} sebagai CSS. Aplikasi ini tidak mengambil foto dan lokasi saat melakukan pencatatan presensi dan tidak terdapat sistem untuk mendeteksi dan menghitung jam lembur. Sedangkan aplikasi penelitian skripsi presensi \textit{online} dikembangkan dengan metode \textit{scrum} dan menggunakan \textit{framework Laravel} 9. Aplikasi penelitian skripsi presensi \textit{online} mengambil foto dan lokasi saat pencatatan presensi dan memiliki sistem untuk mendeteksi dan menghitung jam lembur.

Manu dan Benufinit~\cite{manu2020pengembangan}, melakukan penelitian tentang pengembangan sistem absensi \textit{online} berbasis \textit{web} menggunakan \textit{maps javascripts API}. Pengembangan aplikasi ini menggunakan metode RAD (\textit{Rapid Application Development}) dan menggunakan bahasa pemrograman PHP (\textit{Personal Home Page}) yaitu \textit{PHPMaker}. Aplikasi ini tidak mengambil foto saat melakukan presensi dan tidak memiliki sistem untuk mendeteksi dan menghitung jam lembur. Sedangkan aplikasi penelitian skripsi presensi \textit{online} dikembangkan dengan metode \textit{scrum} dan menggunakan \textit{framework Laravel} 9. Aplikasi penelitian skripsi presensi \textit{online} mengambil foto saat melakukan presensi dan terdapat sistem untuk mendeteksi dan menghitung jam lembur. 

Febriandirza~\cite{febriandirza2020perancangan}, melakukan penelitian terkait perancangan aplikasi absensi \textit{online} dengan menggunakan bahasa pemrograman Kotlin. Aplikasi ini dikembangkan dengan menggunakan metode \textit{agile} dan menggunakan bahasa pemrograman \textit{agile} berbasis Android. Aplikasi ini tidak mengambil foto saat presensi dan tidak ada sistem untuk mendeteksi dan menghitung jam lembur. Sedangkan aplikasi penelitian skripsi presensi \textit{online} dikembangkan menggunakan metode \textit{scrum} dan menggunakan bahasa pemrograman PHP (\textit{Personal Home Page}) dengan \textit{framework Laravel} 9. Aplikasi penelitian skripsi presensi \textit{online} mengambil foto saat presensi dan terdapat sistem untuk mendeteksi dan menghitung lembur.

Pesik dan Tanaem~\cite{pesik2022perancangan}, melakukan penelitian terkait perancangan sistem informasi absensi online deteksi lokasi berbasis \textit{web} untuk mempermudah pegawai kantor untuk melakukan absensi. Aplikasi ini dikembangkan menggunakan metode \textit{waterfall} dan menggunakan bahasa pemrograman PHP (\textit{Personal Home Page}) dan \textit{framework Bootstrap} sebagai CSS. Aplikasi ini sama-sama menggunakan fitur pengambilan lokasi dan foto saat presensi tetapi aplikasi ini tidak terdapat sistem untuk mendeteksi dan menghitung lembur. Sedangkan aplikasi penelitian skripsi presensi \textit{online} dikembangkan dengan metode \textit{scrum} dan menggunakan \textit{framework Laravel} 9 dan \textit{framwork Tailwind} sebagai CSS. Aplikasi penelitian skripsi presensi \textit{online} memiliki fitur untuk mendeteksi dan menghitung lembur.

Rumbang dan Tony~\cite{rumbang2024perancangan}, melakukan penelitian terkait perancangan aplikasi presensi berbasis web untuk
karyawan magang di PT. Sembilan Pilar Semesta agar mempermudah karyawan magang untuk melakukan presensi dan membantu pihak manajemen PT. Sembilan Pilar Semesta dalam melakukan manajemen data. Aplikasi ini dikembangkan dengan \textit{framework Laravel} dan metode \textit{waterfall}. Pada aplikasi ini tidak terdapat halaman untuk registrasi akun sehingga \textit{admin} yang harus menambah akun baru. Aplikasi ini mengambil lokasi karyawan untuk menentukan WFO atau WFH dari lokasi yang didapat saat presensi tanpa pengambilan foto. Sedangkan pada penelitian skripsi presensi \textit{online} menggunakan metode \textit{scrum} dengan \textit{framework} yang sama yaitu \textit{Laravel}. Aplikasi ini menggunakan GPS untuk pencatatan lokasi dan pengambilan foto untuk validasi atau bukti presensi. Aplikasi ini dilengkapi sistem untuk deteksi dan menghitung jumlah jam lembur \textit{staff} jika presensi dilakukan diluar jam kerja atau hari libur. Aplikasi ini terdapat halaman registrasi namun, untuk mengaktifkan akun tersebut membutuhkan persetujuan \textit{Manager Human Resources}.

